\section{Introduction}
\label{sec:introduction}
A specification is useful during development only if programmers can obtain
feedback from it. For C, that feedback must cover both functional results and
the memory resources that make an execution safe. Separation-logic tools such
as CN express these requirements in a single contract~\cite{pulte2023cn}.
Fulminate makes such contracts executable by translating them into runtime
checks~\cite{banerjee2025fulminate}. This enables testing before a complete
verification proof is available.

Two obstacles remain when testing proof-oriented specifications. First,
fixing an input need not make its specification executable. A least-root
contract can characterize a result by its relation to \emph{every} nonnegative
integer. The program returns a concrete value, but checking that value still
requires reasoning about an unbounded domain. Rewriting the contract into a
finite monitor is possible for some properties; discovering and justifying
that rewrite is additional specification work. Second, the execution used to
test an annotation must retain the information produced by the program.
A verifier can replace a loop by its invariant or a call by its summary.
A test of those annotations needs the actual visited states, including any
facts that the annotation omits.

We present \tespec{}, a method for \emph{proof-based testing of separation-logic
specifications for C}. A test supplies bindings for selected logical inputs.
\tespec{} executes the program symbolically under those bindings, observes its
states at selected specification boundaries, and checks the resulting
propositions in the verifier's logic. It executes available function bodies
and loop iterations. Invariants and contracts generate obligations on copies
of the reached states, while the execution retains the original states.
Automatic reasoning handles routine obligations; residual propositions can be
completed by an agent or developer and checked by the proof assistant.

This combination makes three kinds of information available to the test
oracle: the program's computed values, its owned memory, and the mathematical
meaning of the original assertion. A small example exposes why all three
matter. Suppose a read function promises to return the value in an owned
integer cell, but its implementation returns the negated value. On input
\texttt{INT\_MIN}, the recorded Fulminate configuration exits successfully,
whereas \tespec{} rejects the arithmetic safety obligation: the positive
negation is outside the signed integer range. For quantified least-root and
divisibility contracts, the recorded instrumenter cannot infer a finite
quantifier bound, whereas proof completion checks the original obligations.
Section~\ref{sec:case-studies} reconstructs these cases and their controls.

The central technical problem is preserving the meaning of a test while
reusing verification infrastructure. Assertions must observe execution without
restricting it. Heap simplification must retain the values later assertions
will inspect. An auxiliary assertion may omit unrelated stack slots, but it
must not silently discard user-owned memory. \tespec{} addresses these problems
with copied-state observations, joint heap/value normalization, and exact
local-storage frames. A nonempty entry domain and complete execution coverage
then connect successful obligation checking to successful instance execution.

Our contributions are:
\begin{enumerate}
\item \textbf{A stateful proof-based testing method.}
Typed input bindings, actual body execution, and copied-state observations
turn original separation-logic assertions into test obligations at function,
loop, and intermediate boundaries (Section~\ref{sec:approach}).
\item \textbf{A formal account of instance checking.}
We define the tested domain and observation semantics, establish reference
trace preservation, and mechanize nonvacuous acceptance under explicit
execution-coverage interfaces. Exact frames and joint normalization state
which memory and value information a checker must preserve
(Section~\ref{sec:semantics}).
\item \textbf{An implementation and comparative case studies.}
The implementation reuses \qcp{}'s symbolic execution and Rocq obligations.
Evidence-backed cases examine an integer-boundary miss, quantified properties,
and spatial proof completion; a separate evaluation protocol specifies the
broader comparisons and their cost units
(Sections~\ref{sec:case-studies} and~\ref{sec:evaluation}).
\end{enumerate}
The testing boundary is explicit: acceptance covers the domain selected by the
input bindings. Full verification checks the original contract for all its
admissible inputs. We compare these workflows directly, rather than treating
a verifier with an added input equality as an execution-preserving test.
